\documentclass[10pt]{article}
\usepackage[margin=0.85in]{geometry}
\usepackage{booktabs,enumitem,hyperref,amsmath}
\setlist{itemsep=1pt,topsep=2pt}
\pagestyle{empty}
\begin{document}
\noindent{\Large\bfseries DP-EAMxx DYCOMS-II RF02: cloud sedimentation $\times$ precip suppression}\\[2pt]
Plan, 6 Nov 2026. Prepared by Claude Code (AI assistant) on behalf of Naser Mahfouz, for Naser and Johannes.

\section*{Scope}
EAMxx only (DP-EAMxx, P3/SHOC/SPA); EAM/MG2 serves only as the reference for the sedimentation knob.
Extends the ACI $\times$ resolution study (scmlib \texttt{aci-res-study}): same case, domains, time steps and Nd targets.
All on pm-gpu. Sub-km runs (625, 312, 156 m) are \textbf{held} until we have looked at these results.

\section*{Matrix (144 runs, 6 h each, 3D turbulence off)}
\begin{tabular}{@{}ll@{}}
\toprule
factor & levels \\ \midrule
dx & 40, 20, 10, 5, 2.5, 1.25 km \\
Nd target (cm$^{-3}$) & 20, 50, 100, 200 (current SPA factors; diagnosed Nd runs 10--25\,\% high) \\
cloud sedimentation factor & 0.5, 1, 2 \\
precip suppression & on (KK2000 default), off (Nd exponent 0, prefactor compensated) \\ \bottomrule
\end{tabular}

\smallskip\noindent The (sed 1, suppression on) cells repeat the earlier ACI-study runs, which serves as a consistency check.
Estimated cost: $\sim$10--20 node-hours, run packed in debug ($\le$8 nodes) and interactive allocations.

\section*{New knobs}
\begin{itemize}
\item \textbf{\texttt{p3::cloud\_sedimentation\_factor}} (new, default 1.0) multiplies P3's mass- and number-weighted
cloud-droplet fall speeds, the counterpart of MG2's \texttt{cld\_sed}, which scales the Stokes prefactor \texttt{acn}.
Droplet fall speeds used for riming are unchanged.
Branch \texttt{mahf708/eamxx/p3-cloud-sed-factor} on top of E3SM PR \#9017.
\item \textbf{Precip suppression off}: P3 autoconversion is $1350\,q_c^{2.47}N_d^{-1.79}$ ($N_d$ in cm$^{-3}$).
We set the exponent to 0 and the prefactor to $1350\cdot N_{\rm ref}^{-1.79}$ with $N_{\rm ref}=63$\,cm$^{-3}$
(the geometric mean of 20--200), giving \textbf{0.812}. Other reference values: 1.04 at 55\,cm$^{-3}$
(the default SPA Nd at DYCOMS), 1.23 at 50, 0.355 at 100.
Nd still acts through sedimentation, evaporation and LW (SW is off in this nocturnal case).
\end{itemize}

\section*{New output}
\begin{itemize}
\item \textbf{AeroCom cloud-top fields} (PR \#9017 \texttt{aerocom(..)} operator), every 15 min, per column:
cloud-top fraction (total and liquid), $T$, $p$, $r_{\rm eff}$ and in-cloud CDNC at liquid cloud top.
\item \textbf{LWP spatial variance}, computed online every 15 min as the domain mean of
(LWP $-$ domain-mean LWP)$^2$, using the expression syntax (\texttt{mean('col')}, \texttt{broadcast\_like}).
The SD is its square root, cross-checked against the 15-min 2D LWP output.
\end{itemize}

\section*{Steps}
\begin{enumerate}
\item One build (SICE compset, resolution-independent executable); 40 km smoke test of the new knob and output.
\item Run 40 km--2.5 km, then 1.25 km; check Nd, stability and the knob's effect on rain and $w_e$.
\item Analysis as in the ACI study (h4--6 means, susceptibility $d\ln{\rm LWP}/d\ln N_d$, $w_e$, rain, cloud cover,
LWP SD, cloud-top CDNC/$r_{\rm eff}$); results emailed to both of us.
\item Sub-km only after we discuss.
\end{enumerate}
\end{document}
